% =============================================================================
% Related Works section for the H2 paper.
% Standard packages only (amsmath, cite/natbib).
% =============================================================================

\section{Related Work}
\label{sec:related}

Post-training quantization (PTQ) has a long history in neural network
deployment, and the literature divides naturally into methods that use
calibration data to refine the quantized weights, methods that do not,
and methods that preprocess the weights with a structured transformation
before quantization. Our method belongs to the third category, and we
survey the three strands in turn before discussing classical scalar
quantization theory and the specific relationship to the TurboQuant
approach that motivates our construction.

\subsection{Calibration-based post-training quantization}

The dominant paradigm in modern PTQ uses a small calibration dataset of
a few hundred to a few thousand input examples to refine the quantized
weights. BRECQ \citep{li2021brecq} minimizes a block-wise reconstruction
loss under Hessian approximations and produces competitive 4-bit
accuracy on convolutional networks with around one thousand calibration
images. AdaRound \citep{nagel2020adaround} introduces a learnable
rounding policy that is optimized against a reconstruction objective,
and shows that rounding down as opposed to the nearest quantization grid
point can be beneficial at low bit widths. QDrop \citep{wei2022qdrop}
randomly drops quantization during calibration, which improves
generalization of the quantized network. HAWQ \citep{dong2019hawq}
assigns bit widths to different layers based on Hessian spectra,
effectively solving a mixed-precision allocation problem. In the
language-model setting, GPTQ \citep{frantar2022gptq} extends the
calibration paradigm with an approximate second-order update that
operates layer by layer, while AWQ \citep{lin2023awq} scales weight
channels based on activation statistics to preserve the coordinates
most used by downstream computation. SmoothQuant \citep{xiao2023smoothquant}
and OmniQuant \citep{shao2024omniquant} shift the difficulty of
quantization between weights and activations via per-channel scaling
factors that are either hand-designed or learned.

All of these methods require at least some unlabelled examples from the
target distribution, they require gradient or Hessian computation
through the network, and they typically require hours of compute per
model for the calibration step. Our method bypasses all three
requirements. H2 uses no calibration data, computes no gradients, and
runs to completion in under a minute on a single CPU core for ViT-B/16.
We therefore compare head-to-head only against other training-free
baselines, leaving the question of whether calibration-based refinement
can be layered on top of H2 to future work.

\subsection{Training-free quantization}

The simplest training-free scheme, round-to-nearest with per-channel
absolute-maximum scaling (RTN-absmax), remains a common reference point
in the LLM literature and serves as one of our baselines. Data-free
quantization methods such as DFQ \citep{nagel2019dfq} equalize the
weight magnitudes across channels before RTN quantization to reduce the
effective range each channel must cover, but they do not attempt to
reshape the weight distribution. HQQ \citep{badri2023hqq} solves a
half-quadratic regression for the quantization scale and zero point
without any calibration data and achieves strong results for large
language models at 3 to 4 bits.

Classical CNN PTQ literature also produced training-free or
nearly-training-free variants of BRECQ-style methods and of PACT
\citep{choi2018pact}-based clipping, but these either involve small
amounts of training data or rely on specific architectural assumptions
that do not translate to modern transformer-based architectures.

Our H2 method competes in this training-free tier alongside RTN-absmax
and the training-free variant of QuaRot. The empirical comparison in
Section~\ref{sec:experiments} isolates the contribution of the rotation
step and the codebook choice within this tier, and we position H2 as
complementary to rather than competing with the calibration-based
methods above.

\subsection{Rotation-based quantization}

A recent thread of work applies a structured transformation to the
weights before quantization. The Randomized Hadamard Transform and its
close relatives have long been studied in randomized numerical linear
algebra \citep{ailon2006srht} for their ability to spread signal energy
uniformly across coordinates, and the same property makes them attractive
for quantization because it mitigates the outlier problem that plagues
naive RTN. QuaRot \citep{ashkboos2024quarot} formalizes this idea for
large language models and combines a Hadamard rotation with per-channel
absmax uniform quantization, followed optionally by GPTQ-style
refinement for the weights. Their reported results show that the
rotation step alone closes much of the gap between RTN and the full
calibration-based methods at four bits on Llama-family models.

SpinQuant \citep{liu2024spinquant} replaces the random Hadamard with a
learned rotation optimized end-to-end for the downstream quantization
objective, at the cost of requiring gradient computation during the
optimization. QuIP \citep{chee2023quip} rotates via an incoherence-aware
construction and combines rotation with a specific rounding scheme that
respects the rotated geometry. QuIP\# \citep{tseng2024quipsharp} extends
QuIP with Hadamard incoherence and a lattice codebook that jointly
quantizes groups of coordinates, improving the compression-accuracy
trade-off at two bits.

Our work differs from these in two ways. First, the prior rotation-based
methods all combine their rotation with either a uniform grid
(QuaRot, SpinQuant) or a lattice-based vector codebook (QuIP\#). None of
them use the scalar Lloyd-Max codebook matched to the Beta density that
a rotation produces on the unit sphere. Second, our primary empirical
contribution is on vision models rather than on language models. To the
best of our knowledge the ViT-B/16 and ConvNeXt-Tiny results we report at
two and four bits are the first rotation-based PTQ numbers on standard
ImageNet vision benchmarks that preserve accuracy to within a few
percentage points at five to eight times compression.

\subsection{Classical scalar quantization theory}

The construction of an optimal scalar quantizer for a given source
density dates to Lloyd's 1957 Bell Labs memo, later published as
\citet{lloyd1982}, and to Max's independent analysis
\citep{max1960}. The Lloyd-Max fixed-point equations describe the
optimal placement of centroids at conditional means of bins and the
optimal placement of decision boundaries at midpoints of adjacent
centroids. Bennett's earlier analysis \citep{bennett1948} and the
Panter-Dite formula \citep{panter1951} give the high-resolution
asymptotic distortion of an optimal scalar quantizer as a functional of
the source density. The standard textbook treatment in \citet{gershogray1992}
consolidates these results and shows that, for smooth unimodal
densities, Bennett's integral $\big(\int f(x)^{1/3}\, dx\big)^3$
determines the leading constant. Convergence of the Lloyd-Max iteration
to the unique global optimum under log-concavity of the source density
was established by \citet{fleischer1964} and extended by
\citet{kieffer1983}, both of which are cited in our methodology.

This classical theory provides the mathematical machinery for our
codebook step. What is novel in our setting is not the Lloyd-Max
construction itself but the observation that the post-rotation
coordinate distribution is known analytically and admits a codebook
that is globally optimal rather than merely a good heuristic.

\subsection{Relationship to TurboQuant}

The immediate precursor to our method is TurboQuant \citep{turboquant},
which uses the same spherical geometry argument to compress unit-norm
embedding vectors for nearest-neighbour retrieval. TurboQuant applies
a random orthogonal rotation to L2-normalized embeddings and quantizes
each coordinate against the Beta Lloyd-Max codebook on the unit sphere.
The original TurboQuant paper establishes the mathematical framework
but focuses on embedding compression in the retrieval context, where the
objects being compressed are inputs to and outputs of a network rather
than parameters of the network itself.

Our contribution is to transfer this mathematical framework from
embedding compression to weight compression. Three observations make
the transfer nontrivial. First, neural network weights are not
unit-norm. Our Step 1 in Section~\ref{sec:method-rotation} handles this
by per-row L2 normalization, which recovers the unit-sphere setting.
Second, standard PyTorch implementations of multi-head attention store
the query, key, and value projections in a single raw parameter outside
any linear module. A naive iteration misses this parameter entirely,
which we handle in Section~\ref{sec:method-mha} with a dedicated pass.
Third, the distributional assumption that makes the Beta codebook
optimal is exact only when the row is uniformly distributed on the
sphere. For trained weights this assumption is approximate, and the
quality of the approximation depends on the layer fan-in. We analyze
this architectural boundary in Section~\ref{sec:method-boundary} and
verify the rule empirically in Section~\ref{sec:experiments}.

\subsection{Summary of positioning}

H2 sits at the intersection of three literatures. It borrows the
rotation primitive from the recent rotation-based quantization thread,
the Lloyd-Max codebook construction from the classical scalar
quantization literature, and the specific Beta-on-sphere distributional
analysis from the TurboQuant line of work. The combination is, to our
knowledge, new. Where QuaRot and SpinQuant gain accuracy through
rotation plus uniform quantization, H2 gains additional accuracy
through rotation plus a distribution-matched codebook. Where
TurboQuant targets embedding compression, H2 targets weight
compression. The empirical gap between H2 and the closest prior method
(QuaRot-RTN) grows with layer fan-in, consistent with the theoretical
prediction from Section~\ref{sec:method-whybeats} that the advantage of
matching the codebook to the source density scales as $1/d$.
